\emph{Facts used about wave-front sets} (standard distribution theory; Hormander, The Analysis of Linear Partial Differential Operators I, Ch. 8, cited from memory): $\mathrm{WF}(D_{1} + D_{2}) \subset \mathrm{WF}(D_{1}) \cup \mathrm{WF}(D_{2})$; $\mathrm{WF}(D)$ is empty over an open set exactly when $D$ is smooth there; the pullback by a diffeomorphism maps wave-front sets by the transposed derivative.

\emph{The vacuum's sets.} From $(x + i0)^{-1} = -i\int_{0}^{\infty}e^{ipx}\,dp$ one has $W_{0}(u, v) = \tfrac{1}{2\pi}\int_{0}^{\infty}e^{ip(u - v)}\,dp$ and $\tilde W_{0} = \delta - W_{0} = \tfrac{1}{2\pi}\int_{-\infty}^{0}e^{ip(u - v)}\,dp$. Each is smooth off the diagonal, and on the diagonal its singular directions are $(\xi, \eta) = (p, -p)$ with $p$ of one sign only: $\mathrm{WF}(W_{0}) = \Gamma_{+}$ and $\mathrm{WF}(\tilde W_{0}) = \Gamma_{-}$ (sign convention fixed by the first). So $\omega_{0}$ satisfies the condition on every interval.

\emph{(i)} Suppose the condition holds on $J$. Then $R = W_{\omega} - W_{0}$ has $\mathrm{WF}(R) \subset \mathrm{WF}(W_{\omega}) \cup \mathrm{WF}(W_{0}) \subset \Gamma_{+}$. The anticommutator is the same for every state, $W + \tilde W = \delta\otimes\id_{K}$, so also $R = -(\tilde W_{\omega} - \tilde W_{0})$ and $\mathrm{WF}(R) \subset \Gamma_{-}$. As $\Gamma_{+} \cap \Gamma_{-} = \emptyset$, $\mathrm{WF}(R)$ is empty over $J \times J$ and $R$ is smooth there. Conversely, if $R$ is smooth then $\mathrm{WF}(W_{\omega}) = \mathrm{WF}(W_{0} + R) = \Gamma_{+}$ and $\mathrm{WF}(\tilde W_{\omega}) = \mathrm{WF}(\tilde W_{0} - R) = \Gamma_{-}$. This is the generator analogue of the bulk argument by which Hadamard states differ by smooth kernels (Radzikowski 1996 for the scalar field; for the Dirac field on both orderings, Hollands 2001, where the conditions on the two orderings are imposed separately and the smooth difference is stated: Fact \claimref{synthesis.imported-microlocal-hadamard-condition}). Parameter independence: a smooth regular reparametrisation $f$ ($f' > 0$) acts at weight $\tfrac12$ by $W \mapsto f'(u)^{1/2}f'(v)^{1/2}W(f(u), f(v))$; the smooth factor does not enlarge wave-front sets, and the pullback by $(u, v) \mapsto (f(u), f(v))$ maps $(u, u; \xi, -\xi)$ to $(u, u; f'(u)\xi, -f'(u)\xi)$, preserving the sign of $\xi$; so $\Gamma_{\pm}$ go to $\Gamma_{\pm}$ and the condition is the same in every smooth regular parameter (consistent with Theorem \claimref{synthesis.local-normality-at-the-edge-gives-the-null-ground-state}(iii), where the vacua of two regular parameters differ by a continuous kernel).

\emph{(ii)} A smooth $R$ is bounded on every compact $I \times I \subset J \times J$, hence square integrable there: the local Hilbert-Schmidt condition holds on every closed bounded subinterval of $J$. The rest is a corollary of Theorem \claimref{synthesis.local-normality-at-the-edge-gives-the-null-ground-state}: by its (ii) the scaling limit at every point of $J$ exists and is $\omega_{0}$, which is the local form of (N$_{+}$) there; nothing in that step uses a symmetry of $\omega$ or a temperature. If $J$ contains the edge and $\omega$ is stationary for the boost (the weight-$\tfrac12$ dilation about the edge, Theorem \claimref{synthesis.stationary-horizon-state-is-its-own-edge-scaling-limit}(i)), its (iv) gives $\omega = \omega_{0}$, which satisfies (N$_{+}$) (Theorem \claimref{synthesis.positive-null-energy-gives-r-premises}(i)). Nothing new is claimed for the stationary clause beyond the stronger hypothesis.

\emph{(iii)(a)} With $y = \pi x/\beta$, $x = u - v$, $\tfrac{2\beta}{i}(W_{\beta} - W_{0}) = \sinh(y)^{-1} - y^{-1} = y^{-1}(g(y^{2}) - 1)$, where $g$ is the inverse of the power series $\sinh(y)/y = \sum_{k}y^{2k}/(2k + 1)!$, which has constant term $1$; the $i0$ prescriptions of the two poles at $y = 0$ coincide and cancel. The check inverts the series exactly to order $y^{18}$ and confirms $g_{n} = 2(1 - 2^{2n - 1})B_{2n}/(2n)!$ with the Bernoulli numbers from their recurrence, $g_{0} = 1$, $g_{1} = -\tfrac16$, $g_{2} = \tfrac{7}{360}$; so the remainder is $y$ times a power series, with leading term $-\tfrac{y}{6}$, that is $R = -\tfrac{i\pi}{12\beta^{2}}x + O(x^{3})$. On the real line every Taylor coefficient of $\sinh(y)/y$ is positive, so $\sinh(y)/y \ge 1$ (checked on partial sums at rational points) and $\sinh(y)^{-1} - y^{-1}$ is real-analytic on all of $\R$; hence $R$ is smooth on the whole plane and, by (i), the condition holds on every interval. In symbols: the occupation $f = (1 + e^{-\beta p})^{-1}$ has $|f - \theta| \le e^{-\beta|p|}$ at both ends (checked at rational $q = e^{-\beta}$ and integer $p$), the translation-invariant form of the two-ordering condition: $f$ decays rapidly as $p \to -\infty$ and $1 - f$ as $p \to +\infty$.

\emph{(iii)(b)} The symbol $f_{t} = t\,\theta + (1 - t)(1 - \theta)$ lies in $[0, 1]$, so it defines a gauge-invariant quasi-free state, invariant under the translations. Its remainder symbol is $(1 - t)(1 - 2\theta) = -(1 - t)\,\mathrm{sign}(p)$, which does not decay at either end, and $W_{t} - W_{0} = (1 - t)(\delta - 2W_{0})$; since $2W_{0} = \tfrac{i}{\pi}\mathrm{PV}(x^{-1}) + \delta$, $R = -(1 - t)\tfrac{i}{\pi}\mathrm{PV}(x^{-1})$. For a translation-invariant kernel $k$ on $[0, L]^{2}$, $\|k\|^{2} = 2\int_{0}^{L}(L - x)|k(x)|^{2}\,dx$; for $k = x^{-1}$ restricted to $x > \varepsilon$ this exceeds $\tfrac{L}{2}\int_{\varepsilon}^{L/2}x^{-2}\,dx = \tfrac{L}{2\varepsilon} - 1$, unbounded as $\varepsilon \to 0$ (checked at $\varepsilon = L/2^{k + 1}$: the bound is $2^{k} - 1$). So the local Hilbert-Schmidt condition, and by (ii) the microlocal one, fail on every interval. For $t = \tfrac12$ this is the tracial state.

\emph{(iii)(c)} By Theorem \claimref{synthesis.stationary-horizon-state-is-its-own-edge-scaling-limit} the kernels of $\omega_{c}^{\times}$ differ from $W_{0}$ by a function analytic off the edge (exact example at $c = \tfrac12$ on a half: $-\tfrac{i}{2\pi}(a - b)/((a + b)(a^{2} + b^{2}))$ at $u = a^{2}$, $v = b^{2}$, a rational function with positive denominator for $a, b > 0$), and on the cross-edge block by $-\tfrac{i}{2\pi}(u - v)^{-1}$, whose value at $(s, -s)$ is $-\tfrac{i}{2\pi}(2s)^{-1}$; both are homogeneous of degree $-1$ about the edge (checked), hence unbounded on every neighbourhood of $(0, 0)$ and not smooth there. By (i) the condition holds on intervals avoiding the edge and fails on intervals containing it. The check also confirms that the boost rotor of $\Cl(1,3)$ dilates the generator about its edge, which turns stationarity into this homogeneity.

\emph{(iii)(d)} The symbol $f = 1 - c/(1 + p)$ for $p > 0$ and $c/(1 + |p|)$ for $p < 0$ lies in $[0, 1]$ for $0 < c \le \tfrac12$. Its remainder symbol $r(p) = -\mathrm{sign}(p)\,c/(1 + |p|)$ is odd and square integrable ($\int r^{2} = 2c^{2}\int_{0}^{\infty}(1 + p)^{-2}dp = 2c^{2}$; the check bounds the integer partial sums and confirms the oddness), so the translation-invariant remainder kernel is square integrable and the local Hilbert-Schmidt condition holds on every bounded interval. By oddness, $R(x) = \tfrac{1}{2\pi}\int e^{ipx}r(p)\,dp = -\tfrac{ic}{\pi}\int_{0}^{\infty}\tfrac{\sin(px)}{1 + p}\,dp$, and for $x > 0$ the substitution $s = px$ gives $\int_{0}^{\infty}\tfrac{\sin s}{x + s}\,ds$, which is bounded by $\int_{0}^{\pi}\tfrac{\sin s}{s}\,ds$ (the lobes alternate in sign with decreasing size) and tends to $\int_{0}^{\infty}\tfrac{\sin s}{s}\,ds = \tfrac{\pi}{2}$ as $x \to 0^{+}$. So $R$ is bounded and $R(x) \to -\tfrac{ic}{2}\,\mathrm{sign}(x)$ as $x \to 0$: a jump across the diagonal, so $R$ is not smooth (not even continuous) there, and by (i) the microlocal condition fails at every point although the remainder is bounded and square integrable. Equivalently, $|p\,r(p)| \ge c/2$ for $|p| \ge 1$ (checked), so $r$ is not rapidly decreasing.

\emph{(iv)} Immediate from (iii)(a) and (ii): the affine-KMS states are not ground states for $\beta < \infty$ (their symbol is not $\theta$), yet satisfy the condition, hence the local Hilbert-Schmidt condition, at every point. The Reading of Theorem \claimref{synthesis.local-normality-at-the-edge-gives-the-null-ground-state} states the second half as argued; (iii)(a) proves it.

\emph{Reading (argued, not part of the statement).} What the outside results bring, and what they assume. Radzikowski's microlocal form of the Hadamard condition, and Hollands' version for the Dirac field on both orderings (Fact \claimref{synthesis.imported-microlocal-hadamard-condition}), are the bulk conditions of which the condition here is the generator analogue. No theorem restricting a four-dimensional Hadamard state to the algebra of one null generator is in the store, and for a four-dimensional Dirac field the restriction to one null line is not defined without transverse smearing; the series reaches one chiral field per light ray only through the s-wave reduction J3 (Kay and Wald, Phys. Rep. 207 (1991) 49, analyse restrictions to a horizon, from memory). Verch's theorem (for the scalar field) and its Dirac analogue (D'Antoni and Hollands), Fact \claimref{synthesis.imported-local-quasi-equivalence-of-hadamard-states}, make local quasi-equivalence of quasi-free Hadamard states a theorem in the bulk, the principle that Fact \claimref{synthesis.imported-haag-local-normality}(b) states; on a generator, a smooth remainder is trace class on bounded intervals, and with the Powers-Stormer inequality $\|A^{1/2} - B^{1/2}\|_{2}^{2} \le \|A - B\|_{1}$ (Fact \claimref{synthesis.imported-powers-stormer-takesaki-disjointness}(c)) the condition here would give local normality itself, not only the local Hilbert-Schmidt condition. The bulk condition is propagated by the field equation from a neighbourhood of one Cauchy surface (Fulling, Sweeny and Wald 1978, Fact \claimref{synthesis.imported-propagation-of-hadamard-form}, for the scalar field; the Dirac analogue is not imported), so for a hole formed by collapse whose state is Hadamard before the collapse, and given a restriction theorem, the local form of (N$_{+}$) at every point of every generator, the future horizon included, would follow from one initial condition rather than from a premise posed at every point. On the eternal hole the condition cannot select the state: by (iv) the global datum is still needed, which is the stationarity of the Kay-Wald uniqueness theorem (Fact \claimref{synthesis.imported-kay-wald-uniqueness-and-no-go}(a)). Kay and Wald also give a no-go result for the extended Schwarzschild-de Sitter field (no stationary Hadamard state; a gap in the Kay-Wald proofs was later filled by Kay and Lupo, App. B; Kerr is the case described as argued in Kay's review): Fact \claimref{synthesis.imported-kay-wald-uniqueness-and-no-go}(b), (c); with a positive cosmological constant the stationary-class route to (N$_{+}$) at a black-hole horizon then has no state to describe, and the stationarity-free local route is the one that applies there.

\emph{Relation to the condition at infinity (\claimref{synthesis.normality-on-the-closed-generator-gives-the-vacuum-folium-not-the-ground-state}).} The condition here is ultraviolet and posed at finite points; normality at the point at infinity of the closed generator is infrared. Neither implies the other: the affine-KMS states, and the vacua of parameters whose affine slopes differ at the two ends with a $C^{\infty}$ transition, satisfy the condition here on every bounded interval and fail the one at infinity; the rank-one excitation $P_{0} + |g\rangle\langle g|$ of the vacuum projection $P_{0}$, with $g$ a unit vector orthogonal to its range and not smooth, satisfies the one at infinity and fails this one; and with $g$ a Schwartz function it satisfies both and is not a ground state, so even together they do not give (N$_{+}$) without stationarity. The slope-jump states and the half-thermal products of (iii)(c) share one obstruction, a remainder homogeneous of degree $-1$, at infinity for the former and at the edge for the latter. Since the condition here holds in every smooth regular parameter, and the closed generator makes infinity a regular point with the spinor weight, the condition imposed on every arc of the closed generator is well posed and implies normality at infinity (argued). Imposed at every point of a local horizon, either condition would also fix the period that the Clausius relation on the action shell leaves free in matter-free regions (\claimref{synthesis.clausius-on-the-action-shell-ties-entropy-density-to-period}): the period-$2\pi c$ model there fails the condition here at the patch edge by (iii)(c). That would be a new premise, not a derivation. The result is transport-absent: it uses only the continuum mode space and the stored local-normality results.
